\documentclass[11pt,a4paper]{article}

% ---------- Packages ----------
\usepackage[utf8]{inputenc}
\usepackage[T1]{fontenc}
\usepackage{lmodern}
\usepackage{microtype}
\usepackage{geometry}
\usepackage{setspace}
\usepackage{hyperref}
\usepackage{booktabs}
\usepackage{enumitem}
\usepackage{titlesec}

\geometry{margin=1in}
\onehalfspacing
\hypersetup{
  colorlinks=true,
  linkcolor=black,
  urlcolor=blue,
  citecolor=black
}

\title{\textbf{Automated Classification of Scientific Conferences via Machine Learning}\\
\large An Evaluation Framework for the Deanship of Research}
\author{}
\date{}

\begin{document}
\maketitle

\begin{abstract}
The rapid expansion of scientific conferences and open-access publishing has introduced
significant administrative and ethical challenges for research governance. Manual validation
processes are increasingly inefficient and inconsistent. This report presents a comprehensive
machine learning framework for automated conference classification designed to support
institutional decision-making within the Deanship of Research. The system integrates semantic
analysis, structured metadata evaluation, ontology-based mapping, and integrity detection
mechanisms. Emphasis is placed on explainability, human oversight, and long-term validation
to ensure transparency, robustness, and alignment with institutional quality standards.
\end{abstract}

\section{Introduction}

The scientific communication ecosystem has undergone structural transformation over the
past decade. Digitization, open-access proliferation, and the emergence of new publishing
models have dramatically increased the number of conferences competing for academic
submissions. Research administrations must now evaluate applications across multiple
dimensions simultaneously: thematic relevance, academic prestige, ethical legitimacy,
and compliance with institutional funding policies.

Traditional manual review processes are time-consuming and susceptible to variability
across evaluators. Furthermore, predatory conferences have adopted increasingly
sophisticated tactics, including name mimicry and artificial prestige signaling.
An automated classification framework provides a scalable infrastructure that standardizes
evaluation criteria while preserving human authority over final decisions.

\section{Conceptual and Methodological Foundations}

Conference classification is best modeled as a multi-label text classification problem.
Each event may simultaneously belong to multiple conceptual categories, including topical
domain, prestige tier, and legitimacy status. Supervised learning algorithms such as
Naive Bayes, Support Vector Machines (SVM), and Random Forest have historically
demonstrated strong performance in high-dimensional textual representations.
However, the semantic complexity of Calls for Papers and abstracts requires deeper
contextual modeling.

Transformer-based architectures such as BERT and DistilBERT offer contextual embeddings
that capture semantic relationships beyond frequency-based approaches such as TF--IDF.
These models enable the system to distinguish subtle differences in scope and detect
semantic inconsistencies indicative of low-quality venues.

\section{System Architecture}

The proposed system is structured as modular infrastructure integrated into a Research
Management System (RMS). It consists of four principal layers: ingestion, semantic analysis,
feature engineering, and integrity evaluation.

The ingestion layer extracts metadata including title, abstract, keywords, DOI, acronym,
and URL. Identifier normalization and cross-referencing with authoritative databases
(e.g., DBLP, Crossref, CORE) reduce ambiguity and ensure persistent traceability.

Semantic classification combines embedding-based similarity measures with ontology-guided
mapping. By projecting submissions and venue descriptions into shared vector spaces,
cosine similarity scores quantify thematic alignment. Ontology-driven classification,
such as hierarchical taxonomies, enhances interpretability and consistency in reporting.

Feature engineering operationalizes institutional standards. Academic rigor is approximated
through signals such as acceptance rates, program committee reputation, historical continuity,
and bibliometric proxies. Thematic consistency is evaluated via embedding similarity between
the submission and prior proceedings. Transparency indicators include clarity of peer-review
procedures, affiliation with recognized societies, and stability of digital identifiers.

\section{Predatory Conference Detection}

A dedicated integrity component classifies venues according to legitimacy risk. Machine-readable
indicators include hyper-broad scope entropy, aggressive solicitation patterns, name similarity
to established conferences, opaque committee structures, and abnormal publication volume spikes.
String distance metrics, Named Entity Recognition, and anomaly detection algorithms support
binary or probabilistic risk scoring.

\section{Model Portfolio and Selection Strategy}

Given institutional requirements for both performance and interpretability, a hybrid
model portfolio is recommended in the Literature. As shown in Table~\ref{tab:modelportfolio}, each model
contributes complementary strengths within the classification pipeline.

\begin{table}[h]
\centering
\caption{Recommended Model Portfolio for Institutional Conference Classification}
\label{tab:modelportfolio}
\begin{tabular}{@{}p{2.5cm}p{3.2cm}p{7.5cm}@{}}
\toprule
\textbf{Model} & \textbf{Primary Role} & \textbf{Institutional Strength} \\
\midrule
Naive Bayes & Rapid thematic labeling & Computational efficiency for large-scale metadata processing. \\
SVM & Multi-label quality classification & Robustness in high-dimensional sparse text spaces. \\
Random Forest & Standards verification & High interpretability and feature importance reporting. \\
CNN & CFP pattern detection & Effective recognition of adversarial textual templates. \\
BERT / DistilBERT & Deep semantic analysis & Context-aware language understanding for legitimacy inference. \\
XGBoost & Structured metadata scoring & High precision in binary classification tasks. \\
\bottomrule
\end{tabular}
\end{table}

The integration of these models into an ensemble architecture enables semantic depth
while preserving structured decision transparency.

\section{Evaluation Framework}

Classifier performance must be evaluated through accuracy, precision, recall,
and F1 score. Precision is critical to avoid approving low-quality venues,
while recall ensures legitimate conferences are not unjustly rejected.
Temporal validation procedures are necessary to guarantee model stability
across evolving publication ecosystems. Grouped validation by discipline
and region detects structural biases. Class imbalance mitigation techniques,
including weighted loss functions and threshold calibration, are required
to prevent skewed decision boundaries.

\section{Governance, Explainability, and Equity}

Automation bias represents a significant institutional risk. The framework must
enforce a human-in-the-loop structure in which automated outputs serve as
decision support rather than final authority. Explainable AI mechanisms,
including feature attribution scores and structured reason codes, provide
defensible transparency for administrative decisions.

Equity considerations require diversification of quality indicators beyond
commercial rankings. Emerging research communities must not be systematically
disadvantaged due to limited bibliometric visibility. Periodic bias audits
and retraining cycles are therefore integral components of governance.

\section{Implementation Roadmap}

Implementation within the Deanship of Research should proceed through the
definition of institutional standards and the construction of a labeled
gold-standard dataset. Following ingestion pipeline deployment and identifier
normalization, hybrid model training can be executed. Integration into the
RMS enables automated alert generation, structured reporting, and escalation
to human reviewers. Continuous monitoring, retraining, and annual governance
audits ensure long-term robustness.

\section{Conclusion}

Automated conference classification constitutes a foundational infrastructure
for modern research governance. By combining semantic modeling, structured
metadata analysis, and integrity detection within a transparent and accountable
framework, the Deanship of Research can achieve scalable, consistent, and
ethically robust evaluation processes. When coupled with human oversight,
this approach strengthens institutional reputation while protecting researchers
from engagement with predatory venues.

\section*{References}

\begin{enumerate}[leftmargin=1.25em]
\item A Systematic Literature Review on Machine Learning for Automated Requirements Classification - IEEE Xplore \url{https://ieeexplore.ieee.org/document/9307778/}
\item API Learning Hub - Crossref \url{https://www.crossref.org/learning/}
\item Frequently Asked Questions - dblp,  \url{https://dblp.org/faqcont}
\item A Systematic Literature Review on Multi-Label Classification based on Machine Learning Algorithms - ResearchGate,  \url{https://www.researchgate.net/publication/361051294_A_Systematic_Literature_Review_on_Multi-Label_Classification_based_on_Machine_Learning_Algorithms}
\item Intelligent Framework for Detecting Predatory publishing Venues - IEEE Xplore,  \url{https://ieeexplore.ieee.org/iel7/6287639/6514899/10056132.pdf}
\item The CSO Classifier: Ontology-Driven Detection of Research Topics, \url{https://cso.kmi.open.ac.uk/cso-classifier/}
\item CORE Rankings Portal - core.edu.au,  \url{https://www.core.edu.au/conference-portal}
\item PreDefense: Defending Underserved AI Students and Researchers from Predatory Conferences \url{https://proceedings.mlr.press/v142/chen21a/chen21a.pdf}
\item sustainet-guardian/aletheia-probe - GitHub \url{https://github.com/sustainet-guardian/aletheia-probe}
\item An open automation system for predatory journal detection - PMC \url{https://pmc.ncbi.nlm.nih.gov/articles/PMC9940686/}
\item TechDispatch \#2/2025 - Human Oversight of Automated Decision-Making, \url{https://www.edps.europa.eu/data-protection/our-work/publications/techdispatch/2025-09-23-techdispatch-22025-human-oversight-automated-making_en}
\end{enumerate}

\end{document}
